\begin{problem}[Reciprocal Sum Inequality via AM-HM]
Let $a, b, c$ be positive real numbers.
\begin{enumerate}
    \item[(i)] Prove the AM-HM inequality: $\frac{a+b+c}{3} \ge \frac{3}{\frac{1}{a} + \frac{1}{b} + \frac{1}{c}}$.
    \item[(ii)] Hence show that $(a+b+c)\left(\frac{1}{a+b} + \frac{1}{b+c} + \frac{1}{c+a}\right) \ge \frac{9}{2}$.
\end{enumerate}
\end{problem}

\begin{hint}
\rotatebox{180}{\parbox{0.9\textwidth}{
For (i): Apply AM-GM to $\{a,b,c\}$ and $\{\frac{1}{a},\frac{1}{b},\frac{1}{c}\}$, then multiply. For (ii): Use the form $(x+y+z)\left(\frac{1}{x}+\frac{1}{y}+\frac{1}{z}\right) \ge 9$ with substitution $x=a+b$, $y=b+c$, $z=c+a$.
}}
\end{hint}

\begin{solution}
\textbf{Part (i):} Apply AM-GM to $a, b, c$:
\[ a+b+c \ge 3\sqrt[3]{abc} \]

Apply AM-GM to $\frac{1}{a}, \frac{1}{b}, \frac{1}{c}$:
\[ \frac{1}{a} + \frac{1}{b} + \frac{1}{c} \ge 3\sqrt[3]{\frac{1}{abc}} \]

Multiply these inequalities:
\begin{align*}
(a+b+c)\left(\frac{1}{a} + \frac{1}{b} + \frac{1}{c}\right) &\ge 9\sqrt[3]{abc} \cdot \sqrt[3]{\frac{1}{abc}} \\
&= 9
\end{align*}

Divide by $3\left(\frac{1}{a} + \frac{1}{b} + \frac{1}{c}\right)$ to obtain:
\[ \frac{a+b+c}{3} \ge \frac{3}{\frac{1}{a} + \frac{1}{b} + \frac{1}{c}} \]

\textbf{Part (ii):} From (i), we have $(x+y+z)\left(\frac{1}{x} + \frac{1}{y} + \frac{1}{z}\right) \ge 9$ for positive $x, y, z$.

Let $x = a+b$, $y = b+c$, $z = c+a$. Then:
\[ ((a+b) + (b+c) + (c+a)) \left( \frac{1}{a+b} + \frac{1}{b+c} + \frac{1}{c+a} \right) \ge 9 \]

Simplify the left factor:
\[ 2(a+b+c) \left( \frac{1}{a+b} + \frac{1}{b+c} + \frac{1}{c+a} \right) \ge 9 \]

Divide by 2:
\[ (a+b+c) \left( \frac{1}{a+b} + \frac{1}{b+c} + \frac{1}{c+a} \right) \ge \frac{9}{2} \]
\end{solution}

\begin{takeaways}
\item The AM-HM inequality follows from applying AM-GM to both a set and its reciprocals
\item The inequality $(x+y+z)(\frac{1}{x}+\frac{1}{y}+\frac{1}{z}) \ge 9$ is fundamental and versatile
\item Strategic substitution (e.g., $x=a+b$) transforms complex expressions into standard forms
\item Equality holds when $a=b=c$, giving $3a \cdot \frac{3}{2a} = \frac{9}{2}$
\item AM, GM, and HM inequalities: Arithmetic Mean is always greater than or equal to Geometric Mean, which is always greater than or equal to Harmonic Mean. The equalities hold when all numbers are equal.
\item For non-negative real numbers $x_1, x_2, \ldots, x_n$, the power mean of order $p$ is defined as:
\[M_p = \left( \frac{x_1^p + x_2^p + \cdots + x_n^p}{n} \right)^{1/p}\]
For $p > q$, we have $M_p \geq M_q$. This generalizes the AM-GM-HM inequalities.
\end{takeaways}
